\section{Introduction}\label{introduction}

\subsection{Background and Motivation}\label{background-and-motivation}

The proliferation of large language model (LLM) based autonomous agents has intensified interest in multi-agent systems capable of collaborative discovery and knowledge synthesis \cite{ref1}\cite{ref2}. Contemporary frameworks such as LangChain, LlamaIndex, and MetaGPT enable agents to decompose complex tasks, invoke external tools, and maintain conversational state across extended interactions. When multiple such agents operate concurrently---each generating hypotheses, evaluating evidence, and proposing conceptual structures---their outputs accumulate as unstructured natural language prose. This unconstrained textual output, while expressive in semantic range, lacks typed relations, explicit lifecycle tracking, and machine-checkable provenance. The resulting conceptual artifacts are difficult to validate, version, or compose into shared ontological commitments \cite{ref3}\cite{ref4}. Cross-references between concepts remain implicit in natural language, temporal evolution is undocumented, and the absence of a shared schema prevents agents from building upon one another's discoveries in a structured fashion. Each agent may develop its own vocabulary, classification criteria, and confidence calibration, leading to semantic fragmentation that undermines collective intelligence.

Traditional ontology engineering addresses this problem through formal specification: domain experts construct object hierarchies, define property constraints, and establish inference rules using languages such as the Web Ontology Language (OWL) \cite{ref5}. Palantir's Foundry Data Engine (FDE) exemplifies the industrial state of the art, organizing enterprise knowledge around Object Types, Link Types, Property Types, and Action Types configured through a dedicated user interface \cite{ref6}. This approach produces rigorously structured ontologies but imposes substantial deployment friction. Ontology construction requires specialized expertise in both the target domain and the formalism, dedicated tooling that may carry significant licensing costs, and manual curation workflows that do not scale to the velocity of agent-generated conceptual output \cite{ref7}. A further limitation is that LLM-based agents cannot meaningfully participate in this authoring process: they lack programmatic access to ontology construction interfaces and cannot manipulate the underlying formalisms through standard software development kits \cite{ref8}. The Ontology Manager UI in Palantir FDE, for instance, requires human operators to define object types, link types, and action types through point-and-click configuration; no programmatic creation APIs are exposed for automated agent authorship \cite{ref8}. Semantic elements (Object Type, Property Type, Link Type) and kinetic elements (Action Type, Function) must be co-defined through this interface, creating a bottleneck that prevents LLM agents from autonomously extending the ontology as they discover new concepts.

The philosophical foundations of this work draw upon Ludwig Wittgenstein's \emph{Tractatus Logico-Philosophicus} \cite{ref9}. In proposition §1.1, Wittgenstein asserts: \emph{``The world is the totality of facts, not of things.''} The original German---\emph{``Die Welt ist die Gesamtheit der Tatsachen, nicht der Dinge''}---emphasizes that reality consists not of isolated objects but of states of affairs in which those objects participate \cite{ref10}. An object, on this view, has no ontological standing independent of the facts in which it appears. A thing gains its significance only through participation in states of affairs that constitute the fabric of the world. The concept of a validated scientific hypothesis, for instance, is not a static entity possessing an intrinsic property of ``validatedness''; rather, it is constituted by the totality of validation events, evidentiary assertions, and review actions that have occurred in its history. This insight has direct architectural implications: if concepts are understood not as self-subsistent entities but as loci of event histories, then an ontology should be organized around \emph{actions} rather than \emph{objects}. Status, confidence, and relational embeddings become derivative properties computed from the sequence of events that constitute a concept's history, not mutable fields stored alongside the concept itself \cite{ref11}\cite{ref12}. The event-first stance does not deny the existence of objects---Wittgenstein explicitly acknowledges that objects are constituents of facts---but it denies that objects have ontological priority. In engineering terms, this means that the action registry and event schema are designed first, and object-like abstractions emerge as projections of event patterns rather than as foundational data structures. The event-first architecture thus operationalizes Wittgenstein's insight: the ontology records what has happened (the totality of facts), not what purportedly exists (a catalog of things), and concepts acquire their identity through the events in which they participate.

\subsection{Research Gap and Problem Statement}\label{research-gap-and-problem-statement}

Existing operational ontologies, including Palantir FDE and OWL-based knowledge graphs, are predominantly \emph{object-first}: an Object Type defines what entities exist, Property Type describes their attributes, Link Type connects them, and Action Type is added as an afterthought to model operations upon the previously defined object structures \cite{ref6}\cite{ref14}. This design creates deployment bottlenecks at multiple stages of the ontology lifecycle. Concept instances carry mutable status fields that can diverge from their actual event histories when updates occur outside the governed action pathway. Updates are destructive: modifying a concept's properties overwrites prior state, making audit reconstruction dependent on external logging infrastructure that may itself be incomplete or inconsistent. Agents wishing to participate in ontology evolution must navigate complex user interfaces or specialized SDKs with steep learning curves, precluding lightweight automation at the scale required by modern multi-agent research pipelines \cite{ref8}. The cumulative effect is that ontology engineering remains a human-intensive activity disconnected from the generative processes that most need semantic structure---the autonomous production of conceptual content by cooperating agents.

The specific problem addressed in this paper is formulated as follows: \emph{How can an operational ontology achieve production-grade lifecycle traceability, cryptographic integrity, and multi-agent consensus without requiring heavyweight tooling, expert curation, or object-first mutable state?} This problem statement motivates the design of a system that is simultaneously (i) pip-installable and Git-backed, enabling frictionless deployment in standard Python environments without enterprise infrastructure; (ii) natively authorable by LLM agents through Markdown files that serve as both human-readable documentation and machine-parseable semantic assertions, eliminating the need for specialized ontology authoring interfaces; (iii) cryptographically auditable via append-only event chains that detect accidental tampering with the ontological record through SHA-256 hash verification---providing tamper-\emph{evidence} suitable for collaborative audit among trusted participants, though not tamper-\emph{proof} guarantees against adversarial actors without additional cryptographic signing (see Section 3.7); and (iv) semantically governed by a closed action registry with structured precondition validation that eliminates runtime code injection through statically typed comparator enumeration. The system targets the gap between unstructured agent output and formal ontology deployment, where lightweight operational semantics can accelerate conceptual convergence without claiming to replace enterprise platforms for high-assurance compliance workloads \cite{ref13}.

It is important to clarify the security posture that the current system achieves. ADL Lite employs SHA-256 hash chaining to link successive events: each event's hash covers the previous event's identifier, creating a cryptographic chain in which any modification to historical data is detectable as a hash mismatch at the point of first alteration. This mechanism provides \emph{tamper-evidence}---it enables participants to detect accidental or malicious modification of the event record---rather than \emph{tamper-proofness}, which would require digital signatures, a public-key infrastructure, or distributed consensus to prevent alteration by adversarial actors. The current threat model assumes collaborative audit among trusted participants, with the hash chain serving to detect accidental corruption or misattribution. Section 3.7 elaborates the trust model, and Section 3.8 discusses the path to stronger guarantees through Linked Data Proofs \cite{ref21}.

A further deliberate simplification in the current design concerns actor identity. The \texttt{actor} field in each event is a plain string identifier (e.g., \texttt{"discoverer\_agent\_1"}, \texttt{"human\_reviewer"}), not a cryptographically authenticated principal. This means that event authorship is \emph{attributed} but not \emph{proven}: any participant with write access to the shared Git repository can append an event claiming any actor name. This is a deliberate Phase 1 simplification that prioritizes deployment frictionlessness over adversarial resistance. Binding actor identity to cryptographic keys---through Linked Data Proofs with URDNA2015 canonicalization and ed25519 signatures, or via integration with an existing PKI---remains future work (Section 3.7). In the intended deployment context (small-to-medium research teams collaborating within a trusted Git repository), the combination of Git's native audit log (\texttt{git\ blame}, signed commits) and the hash chain's tamper-evidence provides a usable level of accountability.

\subsection{Contributions}\label{contributions}

This paper presents ADL Lite, an event-first operational ontology for multi-agent concept consensus. The system represents concepts as append-only, cryptographically hashed EventChains, with lifecycle state derived exclusively from event history. It is evaluated through eleven experiments (E1--E11) testing chain integrity, status derivation, snapshot round-trip consistency, precondition enforcement, multi-agent auditability, real-time pattern detection, edge offline synchronization, Git baseline comparison, and throughput scalability on real-world financial transaction data. The four principal contributions are enumerated below.

\textbf{Contribution 1: Event-first ontology architecture.} This paper proposes EventChain as the fundamental representational unit for ontological concepts. Each concept is modeled not as a mutable object with stored properties but as an append-only, cryptographically hashed sequence of typed events. A concept's status, confidence score, and validator identity are \emph{computed} from the event chain rather than stored as mutable fields. SHA-256 hashing and cryptographic linking via \texttt{previous\_event\_id} ensure that any tampering with the chain is detectable through \texttt{verify\_integrity()}. The event model includes nine event types spanning the lifecycle (REGISTER, VALIDATE, DEPRECATE, FORK, ARCHIVE), assertions (RELATE, EVIDENCE, SEAL), and communication (ANNOUNCE, PUBLISH). Empirical evaluation on 495,671 EventChains derived from the IBM Anti-Money Laundering (AML) HI-Small dataset demonstrates zero integrity failures across 5,080,714 events processed in 238 seconds (E6), confirming that the architecture scales to production-level data volumes without degradation of cryptographic guarantees.

\textbf{Contribution 2: L4 action blocks with structured precondition validation.} This paper introduces a closed action registry comprising nine core actions (register, validate, fork, deprecate, archive, announce, publish, sync\_dashboard, listen), each defined in \texttt{adl\_core\_ontology.yaml} with declarative precondition rules. Precondition validation employs a \texttt{Comparator} enumeration (EQ, NEQ, GT, GTE, LT, LTE, IN, EXISTS) evaluated against \texttt{ADLFrontMatter} fields through statically typed Pydantic models---eliminating runtime \texttt{eval()} entirely. The \texttt{ActionExecutor} validates preconditions at parse time before dispatching side effects, ensuring that no action can transition a concept into an invalid state. Thirteen test cases covering all registered actions achieve perfect precision (P = 1.0), perfect recall (R = 1.0), and perfect F1 score (F1 = 1.0), with zero false positives and zero false negatives (E4).

\textbf{Contribution 3: Throughput validation and correctness at scale.} This paper reports the results of eleven experiments establishing correctness properties across the full system: 100\% chain integrity verification on 50 valid and 10 intentionally corrupted chains (E1); 100\% status derivation accuracy across 2,204 exhaustive event sequences including 2,197 three-event combinations and 7 edge cases (E2); 100\% snapshot round-trip consistency across 38 concept files spanning synthetic and real AML domains (E3); 100\% precondition enforcement precision and recall on 13 action-validation test cases (E4); 100\% multi-agent auditability across 5-agent simulations producing independently verifiable chains (E5); 100\% chain integrity at scale on 201 chains from real IBM AML transaction data (E6); real-time pattern detection on actual AML transactions with zero false positives across 9,000 legitimate events (E7); and full FDE pipeline validation with 201/201 concepts passing executor validation (E10). These experiments validate system correctness and throughput scalability; they do not constitute domain-level evaluation of ontological quality (e.g., human expert ratings of concept usefulness) or operational compliance certification. The aggregate evidence supports the claim that event-first derivation of ontological state is computationally feasible, cryptographically sound, and empirically correct under exhaustive test conditions.

\textbf{Contribution 4: Open-source reference implementation.} This paper accompanies the release of \texttt{adl\_lite}, a Python toolkit providing an \texttt{ADLParser} for L1--L4 document parsing with Pydantic validation, an \texttt{ActionExecutor} for precondition-validated action dispatch, a \texttt{ConsensusEngine} for status transition enforcement with fork semantics, a \texttt{DataImporter} for bulk EventChain construction from CSV datasets, and a unified experiment runner enabling reproducible evaluation via a single command-line interface. The toolkit is pip-installable, operates over Markdown files stored in standard Git repositories, and requires no external triple store, OWL reasoner, or enterprise infrastructure. All experiments reported in this paper are executable through the provided experiment runner.

\subsection{Paper Organization}\label{paper-organization}

The remainder of this paper is organized as follows. Section 2 surveys related work in ontology engineering, event sourcing patterns, and multi-agent knowledge representation, positioning ADL Lite with respect to Palantir FDE, OWL-based systems, and emerging lightweight ontology approaches. Section 3 presents the ADL Lite architecture in detail, beginning with the philosophical foundation of event-first ontology design, proceeding through the L1--L4 document model, the EventChain data structure, the action executor, the consensus engine, and concluding with the concurrency model, trust model, and semantic web interoperability mappings. Section 4 describes the experimental design across eleven experiments (E1--E11), including the IBM AML dataset and the methodology for integrity, derivation, round-trip, precondition, auditability, and scale evaluation. Section 5 reports quantitative results for all experiments. Section 6 discusses architectural implications, compares ADL Lite with Palantir FDE and OWL-based alternatives, and acknowledges limitations including synthetic pattern detection and stub side effects. Section 7 concludes and identifies directions for future work, including full ConsensusEngine--ActionExecutor integration, cryptographic actor authentication, and expanded dataset evaluation.
